\documentclass[11pt,a4paper]{letter}

\usepackage[T1]{fontenc}
\usepackage[utf8]{inputenc}
\usepackage{lmodern}
\usepackage[margin=2.5cm]{geometry}
\usepackage{microtype}
\usepackage{hyperref}

\signature{Almaz Sapargali \\
           Al-Farabi Kazakh National University \\
           Faculty of Information Technology and Artificial Intelligence \\
           Almaty, Kazakhstan \\
           \texttt{almaz.sapargali2001@gmail.com}}

\address{Almaz Sapargali \\
         Faculty of Information Technology and AI \\
         Al-Farabi Kazakh National University \\
         71 al-Farabi Avenue, Almaty 050040, Kazakhstan \\
         \texttt{almaz.sapargali2001@gmail.com}}

\date{\today}

\begin{document}

\begin{letter}{Editor-in-Chief \\
               \textit{Results in Engineering} \\
               Elsevier}

\opening{Dear Editor,}

We are pleased to submit our manuscript entitled
\textbf{``When Predictive Surrogates Fail as RL Environments: A
Calibrated Physical Twin as Soft Regularizer for HVAC Control''}
for consideration as a research article in \textit{Results in
Engineering}.

\paragraph{Why this work belongs in \textit{Results in Engineering}.}
The journal's scope explicitly includes reproducible
engineering-software contributions in energy systems and building
control, evaluated on standardised benchmarks. Our work addresses a
real obstacle in deploying reinforcement-learning (RL) controllers to
heating, ventilation and air-conditioning (HVAC) systems: standard
RL training pipelines depend on neural-network surrogates whose
predictive accuracy is routinely reported as a quality indicator, but
whose usefulness as closed-loop training environments is rarely
verified. We show that the two properties can decouple sharply, and
we provide a reproducible recipe that resolves the resulting failure
mode. All experiments use the open BOPTEST testbed; all code,
configuration, trained models, and the eleven supplementary
compliance tables (S1--S11) are released alongside the manuscript.

\paragraph{Novelty in four points.}
\begin{enumerate}
  \item \textbf{An explicit negative result with mechanistic
        explanation.} A physically calibrated RC-NeuralODE twin with
        identifiable zone thermal capacitance ($C_\mathrm{zon}$)
        achieves a $0.64\,^{\circ}\mathrm{C}$ 24-hour rollout RMSE
        (versus $1.47\,^{\circ}\mathrm{C}$ uncalibrated) yet
        \emph{collapses} as a stand-alone RL training environment,
        producing live closed-loop RMSE above
        $4\,^{\circ}\mathrm{C}$. We trace this to a structural
        difference between predictive and closed-loop usage: residual
        errors that average out under one-step prediction accumulate
        into a positive-feedback loop under closed-loop training.
  \item \textbf{A hybrid resolution.} We repurpose the calibrated
        twin as a frozen \emph{soft physical regulariser} that enters
        the policy training loss but never the rollout forward pass.
        The resulting backend sustains 1{,}787 environment steps per
        second on a single CPU thread --- an $85\times$ speed-up over
        the standard BOPTEST RTE HTTP--Docker deployment --- and
        restores live closed-loop RMSE below
        $0.65\,^{\circ}\mathrm{C}$.
  \item \textbf{A controller-family-specific finding.} The optimal
        regularisation strength is not universal: $\lambda_\mathrm{temp}=0.10$
        for thermostatic PPO but $\lambda_\mathrm{temp}=0.00$ for
        hierarchical RL and 17-dimensional MORL, with each value
        explained by observation-geometry and policy-hierarchy
        mechanisms. We report this as a falsification of the implicit
        ``one $\lambda$ fits all'' assumption found in much of the
        prior literature.
  \item \textbf{Pre-registered transferability decomposition.}
        Block~3 of the study cross-validates the recipe on three
        BOPTEST hydronic testcases under a manifest committed
        \emph{before} any non-source run. The surrogate component
        transfers structurally (identified $C_\mathrm{zon}$ ratio of
        $1.91 \pm 0.03$ across $N=3$ testcases), while the frozen
        controller, routed through a mechanical hydronic adapter,
        fails on all three --- with comfort-violation failure on the
        residential testcases and energy-inflation failure on the
        commercial one. This component-level decomposition of transfer
        is, to our knowledge, the first of its kind for surrogate-based
        HVAC RL.
\end{enumerate}

\paragraph{Reproducibility and protocol compliance.}
The full pipeline follows the surrogate-development protocol of
Hou \& Evins (2024), with $15$ of the $17$ Reporting-Level-3 items
documented quantitatively in Supplementary Tables S1--S11. Four
git-anchored pre-registration commits (MORL canonical selection,
MORL post-$N{=}5$ falsification, Block~3 pre-registration, Block~3
closure) bound the timeline against post-hoc adjustment.

\paragraph{Suggested reviewers.}
We suggest the following non-conflicted reviewers based on their
recent contributions to RL-HVAC, surrogate modelling, or BOPTEST
benchmarking. The authors have no personal, professional, or
financial relationship with any of them.
\begin{itemize}
  \item Davide Coraci (Politecnico di Torino) --- staged RC
        calibration for building-energy surrogates.
  \item Yan Wang (Carnegie Mellon University) --- integral safety
        metric for HVAC control benchmarking.
  \item David Blum (Lawrence Berkeley National Laboratory) ---
        BOPTEST testbed maintainer and lead author.
\end{itemize}
We also welcome alternative reviewers chosen at the Editor's
discretion.

\paragraph{Originality and conflicts.}
The manuscript has not been published, has not been submitted
elsewhere, and is not under review at any other journal. The author
declares no competing financial interest or personal relationship
that could be perceived as having influenced this work.

We thank you for your consideration and look forward to the reviewers'
comments.

\closing{Sincerely,}

\end{letter}

\end{document}
